\documentclass[../Main.tex]{subfiles}

\begin{document}

\section{Verzeichnisse}

\subsection{Literatur- und Quellenverzeichnis}

\printbibliography[heading=none]

\subsection{Interview Antworten}

\subsection{Testperson 1 Cyrill}
\subsubsection{Pre-Questionnaire}

Das Pre-Questionnaire dient dazu, die Testperson zu charakterisieren. Ihre Aussagen können so im Kontext gesehen und gewichtet werden.

\begin{enumerate}[label=\textbf{P\arabic*}, leftmargin=1.5cm, itemsep=0.5em]
    \item Während der Schulzeit nicht aber in den Ferien etwa die hälfte der Ferien.
    \item Survival und Shooter
    \item Die Grafik, wenn man lange braucht bis man irgendwas versteht und dass man wegrennen kann
    \item Ich stelle mir vor dass man Kämpft und versucht herauszufinden, wie etwas funktioniert.
    \item 3-4
    \item Ganz klar aggressiv und etwas methodisch wenn man das Spiel besser kann
    \item Ds Game kennenlernen, woran ihr jetzt anscheinend so lange gearbeitet habt. Ich habe ja ein paar Ausschnitte bei euch gesehen
    \item Ich finde Unübersichtlichkeit ehrlich gesagt ganz wizig, also es nervt eigentlich richtig, aber es fühlt sich dann belohnend an. Als das schätze ich an einem Game. Zu leichte gegner sind langweilig.
\end{enumerate}

\subsubsection{Post-Questionnaire}

Das Post-Questionnaire erfasst die subjektive Einschätzung der Testperson direkt im Anschluss an den Playtest, pro Designvorschlag sowie für den gesamten Kampfabschnitt. Hinter jeder Aussage steht in Klammern der Designvorschlag, den sie überprüft. Alle Aussagen werden auf einer Skala von 1 (\enquote{trifft gar nicht zu}) bis 5 (\enquote{trifft völlig zu}) bewertet, die offenen Fragen frei beantwortet.

\paragraph{Gegnerverhalten}
\begin{enumerate}[label=\textbf{Q\arabic*}, leftmargin=1.5cm, itemsep=0.5em]
    \item Ich fühlte mich meinen Gegnern grundsätzlich überlegen, ausser in einzelnen zugespitzten Momenten. \hfill (V1)
    \item Auch bei vielen Gegnern gleichzeitig behielt ich die Kontrolle über die Situation. \hfill (V2)
    \item Ich musste bewusst entscheiden, welchen Gegner ich zuerst angreife. \hfill (V3)
    \item Angriffe der Gegner wurden klar angekündigt, bevor sie mich trafen. \hfill (V4)
\end{enumerate}

\paragraph{Gegnerplatzierung und Kampfraum}
\begin{enumerate}[label=\textbf{Q\arabic*}, leftmargin=1.5cm, itemsep=0.5em, resume]
    \item Die Gegner waren zu Beginn klar überschaubar positioniert, das Verhalten wurde danach unvorhersehbarer. \hfill (V5)
    \item Der enge Abschnitt fühlte sich beim Spielen anders an als der offene Abschnitt. \hfill (V6)
    \item Der Kampf hatte spürbare Phasen von Anspannung und Erholung. \hfill (V7)
    \item Neue Gegnertypen wurden mir zuerst einzeln, dann erst in Kombination vorgestellt. \hfill (V8)
\end{enumerate}

\paragraph{Gesamteindruck}
\begin{enumerate}[label=\textbf{Q\arabic*}, leftmargin=1.5cm, itemsep=0.5em, resume]
    \item Insgesamt war dieser Kampfabschnitt für mich befriedigend zu spielen. \hfill (Spielgefühl)
    \item Der gesamte Kampfabschnitt fühlte sich zusammenhängend an, nicht wie eine Aneinanderreihung einzelner Momente. \hfill (Spielfluss)
    \item Der Schwierigkeitsgrad war insgesamt angemessen. \hfill (Schwierigkeit)
    \item Was hat sich am besten angefühlt? \hfill (offen)
    \item Was hat sich unfair oder unangenehm angefühlt? \hfill (offen)
    \item Wie war dieser Kampfabschnitt im Vergleich zu dem, was du dir vor dem Test erwartet hast? \hfill (offen)
\end{enumerate}

\subsubsection{Beobachtungsraster}

Während des Playtests wird das Verhalten der Testperson anhand der folgenden Kriterien dokumentiert, jeweils mit Zeitstempel, um die Beobachtung später dem passenden Segment und dem retrospektiven Interview zuordnen zu können. Zunächst wird nur beobachtet und protokolliert, nicht interpretiert.

\paragraph{Gegnerverhalten}
\begin{enumerate}[label=\textbf{B\arabic*}, leftmargin=1.5cm, itemsep=0.5em]
    \item Verhält sich die Testperson in normalen Phasen risikofreudig und aggressiv (z.\,B. wenig Ausweichen, viel Angreifen) und spürbar vorsichtiger in den zugespitzten Momenten? \hfill (V1)
    \item Verhält sich die Testperson in Momenten mit vielen Gegnern anders (hektischer, mehr Fehler) als in ruhigeren Momenten? \hfill (V2)
    \item Welchen Gegnertyp greift die Testperson zuerst an, wenn Nah- und Fernkämpfer gleichzeitig anwesend sind? \hfill (V3)
    \item Reagiert die Testperson (z.\,B. durch Ausweichen), bevor ein Angriff trifft, oder erst danach bzw. gar nicht? \hfill (V4)
    \item Wiederholt die Testperson denselben Fehler mehrfach an derselben Stelle?
\end{enumerate}

\paragraph{Gegnerplatzierung und Kampfraum}
\begin{enumerate}[label=\textbf{B\arabic*}, leftmargin=1.5cm, itemsep=0.5em, resume]
    \item Zeigt die Testperson zu Beginn eines Encounters ein zielgerichtetes Vorgehen (z.\,B. gezieltes Anvisieren einzelner Gegner) und wird dieses Vorgehen unruhiger, sobald sich die Gegner umzingelnd verhalten? \hfill (V5)
    \item Unterscheidet sich das Bewegungsmuster im engen Abschnitt sichtbar vom offenen Abschnitt (mehr Ausweichen, mehr Distanzhalten)? \hfill (V6)
    \item Zeigt die Testperson in Erholungsphasen ein erkennbar anderes Verhalten (z.\,B. Innehalten, freies Erkunden) als während einer Welle? \hfill (V7)
    \item Wirkt die Testperson beim ersten Auftreten eines neuen Gegnertyps zögerlich oder unsicher, und verschwindet dieses Zögern bei späteren Wiederholungen desselben Typs? \hfill (V8)
\end{enumerate}

\paragraph{Allgemein}
\begin{enumerate}[label=\textbf{B\arabic*}, leftmargin=1.5cm, itemsep=0.5em, resume]
    \item Anzeichen von Verwirrung oder Orientierungslosigkeit (z.\,B. Stehenbleiben, Umherblicken ohne erkennbares Ziel)
    \item Spontane verbale Reaktionen (Wortlaut und Zeitstempel notieren)
    \item Sichtbare Frustration oder sichtbare Freude/Erleichterung, mit Zeitstempel
    \item Auffällige, nicht vorhergesehene Entscheidungen (z.\,B. Ignorieren eines Gegnertyps, unerwartete Fluchtversuche)
    \item Wirkt der Übergang zwischen zwei Segmenten (z.\,B. Welle zu Erholung, eng zu offen) beim Spielverhalten fliessend, oder gibt es einen sichtbaren Bruch (Zögern, Pause, sichtbare Neuorientierung)? \hfill (Spielfluss)
\end{enumerate}

\subsubsection{Semistrukturiertes Interview}

Das Interview findet im Anschluss an das Post-Questionnaire statt und nutzt die Videoaufnahme des Playtests, um gezielt auf die während der Beobachtung notierten Auffälligkeiten einzugehen. Die folgenden Fragen dienen als Leitfaden und werden je nach den tatsächlich beobachteten Auffälligkeiten der jeweiligen Testperson angepasst und ergänzt.

\begin{enumerate}[label=\textbf{I\arabic*}, leftmargin=1.5cm, itemsep=0.5em]
    \item An den markierten Stellen: Was ist dir hier durch den Kopf gegangen?
    \item Du hast hier [Reaktion/Aussage] gezeigt bzw. gesagt. Kannst du beschreiben, warum?
    \item An dieser Stelle bist du wiederholt gestorben bzw. hast du gezögert. Was hat dir das Weiterkommen erschwert?
    \item Vergleiche diesen Moment mit jenem dort: Haben sich die beiden unterschiedlich angefühlt, und wenn ja, wie?
    \item Als hier mehrere Gegner gleichzeitig da waren: Wie hast du entschieden, wen du zuerst angreifst?
    \item Empfandest du diesen Treffer bzw. diesen Tod als eigenen Fehler oder als unfair? Warum?
    \item Gab es einen Punkt, an dem sich der Kampf für dich in einen anderen \enquote{Modus} verschoben hat (z.\,B. von entspannt zu angespannt)? Wo genau?
    \item Gibt es etwas, das du während des Spielens gedacht oder gefühlt hast, das wir bisher nicht angesprochen haben?
\end{enumerate}

\subsection{Testperson 2 Beatrice}

\subsubsection{Pre-Questionnaire}

Das Pre-Questionnaire dient dazu, die Testperson zu charakterisieren. Ihre Aussagen können so im Kontext gesehen und gewichtet werden.

\begin{enumerate}[label=\textbf{P\arabic*}, leftmargin=1.5cm, itemsep=0.8em]
    \item Wie viele Stunden spielst du pro Woche Videospiele? \\
        \textit{Antwort:} Kommt auf die Anzahl Prüfungen an. Unter der woche etwa 3-4h und in den Ferien ca. 12-15h.

    \item Was für Videospielgenres spielst du besonders gerne? \\
        \textit{Antwort:} Shooter, Souls-likes, horror, action, cozy farming sim.

    \item Was macht für dich einen guten Kampfabschnitt aus? \\
        \textit{Antwort:} 1. visuelle coherenz und coole kunst, 2. animierende kämpfe mit vielen verschiedenen moves und kombos, 3. coole waffen, 4. character spezifische Eigenschaften.

    \item Was stellst du dir unter einem Hack and Slash vor? \\
        \textit{Antwort:} Wort wörtlich also dass der Charakter Hackt und Slashed, also seine Gegner

    \item Wie schätzt du dein eigenes Können in actionlastigen Spielen generell ein? (Skala 1--5: \enquote{Anfänger} bis \enquote{sehr erfahren}) \\
        \textit{Antwort:} 3-4

    \item Bevorzugst du eher vorsichtiges, methodisches Spiel oder aggressives, schnelles Spiel? \\
        \textit{Antwort:} aggressiv und schnell

    \item Was erwartest du dir von diesem Playtest, bevor du ihn spielst? \\
        \textit{Antwort:} 1. So die grundrisse eines echten Spiels, also nur das nötigste, 2. Einen spielbaren Charakter mit flüssigen moves, 3. Anfänge eines world building, so man bekommt eine Idee vom konzept, 4. Bossfight und kampfspass.

    \item Worauf reagierst du in Kämpfen erfahrungsgemäss am empfindlichsten: z.\,B. unfaire Treffer, Unübersichtlichkeit, zu leichte Gegner, zu schwere Gegner? \\
        \textit{Antwort:} Zu schwere gegner mit zu punishing moves wie z.b. marlenia.
\end{enumerate}

\subsubsection{Post-Questionnaire}

Das Post-Questionnaire erfasst die subjektive Einschätzung der Testperson direkt im Anschluss an den Playtest, pro Designvorschlag sowie für den gesamten Kampfabschnitt. Hinter jeder Aussage steht in Klammern der Designvorschlag, den sie überprüft. Alle Aussagen werden auf einer Skala von 1 (\enquote{trifft gar nicht zu}) bis 5 (\enquote{trifft völlig zu}) bewertet, die offenen Fragen frei beantwortet.

\paragraph{Gegnerverhalten}
\begin{enumerate}[label=\textbf{Q\arabic*}, leftmargin=1.5cm, itemsep=0.8em]
    \item Ich fühlte mich meinen Gegnern grundsätzlich überlegen, ausser in einzelnen zugespitzten Momenten. \hfill (V1) \\
        \textit{Antwort:} 5
    \item Auch bei vielen Gegnern gleichzeitig behielt ich die Kontrolle über die Situation. \hfill (V2) \\
        \textit{Antwort:} 4
    \item Ich musste bewusst entscheiden, welchen Gegner ich zuerst angreife. \hfill (V3) \\
        \textit{Antwort:} 5
    \item Angriffe der Gegner wurden klar angekündigt, bevor sie mich trafen. \hfill (V4) \\
        \textit{Antwort:} 3
\end{enumerate}

\paragraph{Gegnerplatzierung und Kampfraum}
\begin{enumerate}[label=\textbf{Q\arabic*}, leftmargin=1.5cm, itemsep=0.8em, resume]
    \item Die Gegner waren zu Beginn klar überschaubar positioniert, das Verhalten wurde danach unvorhersehbarer. \hfill (V5) \\
        \textit{Antwort:} 5
    \item Der enge Abschnitt fühlte sich beim Spielen anders an als der offene Abschnitt. \hfill (V6) \\
        \textit{Antwort:} 5
    \item Der Kampf hatte spürbare Phasen von Anspannung und Erholung. \hfill (V7) \\
        \textit{Antwort:} 3
    \item Neue Gegnertypen wurden mir zuerst einzeln, dann erst in Kombination vorgestellt. \hfill (V8) \\
        \textit{Antwort:} 5
\end{enumerate}

\paragraph{Gesamteindruck}
\begin{enumerate}[label=\textbf{Q\arabic*}, leftmargin=1.5cm, itemsep=0.8em, resume]
    \item Insgesamt war dieser Kampfabschnitt für mich befriedigend zu spielen. \hfill (Spielgefühl) \\
        \textit{Antwort:} 4
    \item Der gesamte Kampfabschnitt fühlte sich zusammenhängend an, nicht wie eine Aneinanderreihung einzelner Momente. \hfill (Spielfluss) \\
        \textit{Antwort:} 5
    \item Der Schwierigkeitsgrad war insgesamt angemessen. \hfill (Schwierigkeit) \\
        \textit{Antwort:} 4
    \item Was hat sich am besten angefühlt? \hfill (offen) \\
        \textit{Antwort:} Die Verteilung der Gegner gefiel mir, also wie sie platziert/gespawnt wurden, deren leben war auch angemessen also kein unnötiges leben.
    \item Was hat sich unfair oder unangenehm angefühlt? \hfill (offen) \\
        \textit{Antwort:} es war immer schwere wenn sie in einem haufen gewesen sind. Besonders wenn es die Soliders waren. Springen unangenehm, weil ich erwartet habe das ich so einen krassen hit gegen den boden mach und gegner treffe. Kommt glaube von Elden Ring
    \item Wie war dieser Kampfabschnitt im Vergleich zu dem, was du dir vor dem Test erwartet hast? \hfill (offen) \\
        \textit{Antwort:} Ich habe einen Boss erwartet, welcher nicht gekommen ist. Der Umfang der Spielmechaniken haben mich sehr überrascht. Der Spamgegner also der Solider war angenehm zum spielen. Ich habe mir ein paar mehr sceen elemente vorgestellt, aslo mehr verzierung
\end{enumerate}

\subsubsection{Beobachtungsraster}

Während des Playtests wird das Verhalten der Testperson anhand der folgenden Kriterien dokumentiert, jeweils mit Zeitstempel, um die Beobachtung später dem passenden Segment und dem retrospektiven Interview zuordnen zu können. Zunächst wird nur beobachtet und protokolliert, nicht interpretiert.

\paragraph{Gegnerverhalten}
\begin{enumerate}[label=\textbf{B\arabic*}, leftmargin=1.5cm, itemsep=0.8em]
    \item Verhält sich die Testperson in normalen Phasen risikofreudig und aggressiv (z.\,B. wenig Ausweichen, viel Angreifen) und spürbar vorsichtiger in den zugespitzten Momenten? \hfill (V1)
        \begin{itemize}[nosep]
            \item Die Testperson spielt insgesamt sehr aggressiv, ausser wenn es zu viele Gegner werden (siehe B2).
            \item Beim Heavy Soldier geht sie methodischer vor und versucht, keine Treffer zu nehmen.
            \item Die Heavy Soldiers tankt sie weiterhin einfach. Sie spielt für diesen Gegner zu aggressiv. Vielleicht ist der Wechsel zwischen aggressivem und taktischem Spielstil für sie eine zu starke Umstellung (Vermutung der Beobachtenden).
        \end{itemize}

    \item Verhält sich die Testperson in Momenten mit vielen Gegnern anders (hektischer, mehr Fehler) als in ruhigeren Momenten? \hfill (V2)
        \begin{itemize}[nosep]
            \item Sie sucht Abstand, wenn es zu viele Gegner werden, und fokussiert sich dann neu auf die Fernkämpfer.
            \item Bei zu vielen Gegnern rennt sie weg, um eine Atempause zu machen und zu heilen. Danach geht es aber weiter.
            \item Sie hat sich mittlerweile an Gruppen von vier bis fünf Gegnern gewöhnt.
            \item Sie kann mittlerweile drei Heavy Soldiers gleichzeitig bekämpfen.
        \end{itemize}

    \item Welchen Gegnertyp greift die Testperson zuerst an, wenn Nah- und Fernkämpfer gleichzeitig anwesend sind? \hfill (V3)
        \begin{itemize}[nosep]
            \item Sie versucht meistens zuerst, den Archer auszuschalten. Da aber oft Soldiers im Weg stehen, hat sie hin und wieder Probleme.
            \item Bei niedriger Gesundheit ist ihr primäres Ziel, die Fernkämpfer auszuschalten.
            \item Wird sie von einem Pfeil getroffen, verschiebt sich ihr Fokus auf die Quelle des Schusses.
            \item Sobald sie von irgendeinem Archer getroffen wird, richtet sich ihr Fokus auf alle Fernkämpfer.
        \end{itemize}

    \item Reagiert die Testperson (z.\,B. durch Ausweichen), bevor ein Angriff trifft, oder erst danach bzw. gar nicht? \hfill (V4)
        \begin{itemize}[nosep]
            \item Die Testperson ignoriert eingehende Treffer, hält sie einfach aus und schlägt weiter.
            \item Sie tankt die Treffer weiterhin.
            \item Sie reagiert auf Angriffe erst, nachdem diese sie getroffen haben.
            \item Mittlerweile versucht sie, dem Heavy Archer auszuweichen, aber nur leicht.
        \end{itemize}

    \item Wiederholt die Testperson denselben Fehler mehrfach an derselben Stelle?
        \begin{itemize}[nosep]
            \item Sie konnte seit Beginn nicht richtig springen, weil sie die Leertaste spammt. Das bereitet ihr weiterhin Schwierigkeiten.
            \item Sie versucht wiederholt, sich an einer Kante festzuhalten (Ledge Grab), zielt dabei aber auf eine Wand statt auf die greifbare Box und scheitert deshalb.
            \item Mittlerweile kann sie etwas besser springen.
        \end{itemize}
\end{enumerate}

\paragraph{Gegnerplatzierung und Kampfraum}
\begin{enumerate}[label=\textbf{B\arabic*}, leftmargin=1.5cm, itemsep=0.8em, resume]
    \item Zeigt die Testperson zu Beginn eines Encounters ein zielgerichtetes Vorgehen (z.\,B. gezieltes Anvisieren einzelner Gegner) und wird dieses Vorgehen unruhiger, sobald sich die Gegner umzingelnd verhalten? \hfill (V5)
        \begin{itemize}[nosep]
            \item Bei neuen Encounters geht sie direkt auf den ersten Gegner zu, den sie finden kann.
            \item Zwischen zwei Räumen oder Wellen sucht sie eine Box zum Draufklettern, um den Raum zu überblicken und alle Gegner zu erfassen.
        \end{itemize}

    \item Unterscheidet sich das Bewegungsmuster im engen Abschnitt sichtbar vom offenen Abschnitt (mehr Ausweichen, mehr Distanzhalten)? \hfill (V6)
        \begin{itemize}[nosep]
            \item In offenen Encountern ist sie weniger vorsichtig und geht aggressiver vor. In engeren Encountern scheint sie eher strategisch zu spielen.
            \item In kleinen Räumen verhält sie sich etwas taktischer und weniger aggressiv, möglicherweise weil sie dort weniger Leben hatte. Sie teilt einen Treffer aus und zieht sich wieder zurück, bis alle Gegner tot sind (guerillaartig).
            \item Im kleinen Raum gibt es wenig Platz zum Heilen (siehe B12).
        \end{itemize}

    \item Zeigt die Testperson in Erholungsphasen ein erkennbar anderes Verhalten (z.\,B. Innehalten, freies Erkunden) als während einer Welle? \hfill (V7)
        \begin{itemize}[nosep]
            \item Bevor sie in den nächsten Raum geht, macht sie Pausen, um sich zu erholen und zu heilen.
            \item Zum Heilen geht sie an einen Ort ohne Gegner und nimmt Abstand.
            \item Sie benutzt die Boxen, um daraufzusteigen, kleine Pausen zu machen und zu heilen.
        \end{itemize}

    \item Wirkt die Testperson beim ersten Auftreten eines neuen Gegnertyps zögerlich oder unsicher, und verschwindet dieses Zögern bei späteren Wiederholungen desselben Typs? \hfill (V8)
        \begin{itemize}[nosep]
            \item Sie geht direkt auf die ersten Heavy Soldiers zu und wird hart getroffen.
            \item Beim ersten Treffer durch einen Heavy Soldier ist sie überrascht. Sie wurde nur getroffen, weil sie ihn \enquote{trollen} wollte.
            \item Beim ersten Heavy Archer erkennt sie sofort, dass dessen Geschosse zielsuchend sind und sie verfolgen.
            \item Als der Chain Archer eingeführt wird, ist sie unsicher, ob das alle Gegner im Raum sind, und analysiert den gesamten Raum, bevor sie ihn angreift.
            \item Bei neuen Gegnern hält sie meistens kurz inne, bevor sie angreift.
        \end{itemize}
\end{enumerate}

\paragraph{Allgemein}
\begin{enumerate}[label=\textbf{B\arabic*}, leftmargin=1.5cm, itemsep=0.8em, resume]
    \item Anzeichen von Verwirrung oder Orientierungslosigkeit (z.\,B. Stehenbleiben, Umherblicken ohne erkennbares Ziel)
        \begin{itemize}[nosep]
            \item Nachdem der Encounter im Raum beendet war, bemerkt sie nicht gleich, dass sie weitergehen kann. Auch öffnende Türen bemerkt sie nicht immer sofort.
            \item Wenn sie keinen Gegner sieht, prüft sie die Tür. Ist diese geschlossen, sucht sie den Raum erneut ab.
        \end{itemize}

    \item Spontane verbale Reaktionen (Wortlaut und Zeitstempel notieren)
        \begin{itemize}[nosep]
            \item Keine Einträge notiert.
        \end{itemize}

    \item Sichtbare Frustration oder sichtbare Freude/Erleichterung, mit Zeitstempel
        \begin{itemize}[nosep]
            \item Im kleinen Raum ist sie sichtlich genervt davon, dass sie beim Heilen unterbrochen wird.
            \item Besonders die Heavy Archers scheinen sie zu nerven (Eindruck der Beobachtenden).
        \end{itemize}

    \item Auffällige, nicht vorhergesehene Entscheidungen (z.\,B. Ignorieren eines Gegnertyps, unerwartete Fluchtversuche)
        \begin{itemize}[nosep]
            \item Sie setzt manchmal unnötig Air Combat ein und trifft dann nicht alle Treffer.
            \item Sie verwendet selten Combo-Variationen, überwiegend nur die leichte Attacke.
        \end{itemize}

    \item Wirkt der Übergang zwischen zwei Segmenten (z.\,B. Welle zu Erholung, eng zu offen) beim Spielverhalten fliessend, oder gibt es einen sichtbaren Bruch (Zögern, Pause, sichtbare Neuorientierung)? \hfill (Spielfluss)
        \begin{itemize}[nosep]
            \item Beim Wechsel von Raum zu Raum bemerkt sie nicht immer sofort, dass sich die Tür geöffnet hat (siehe B10). Davor erholt sie sich in der Regel (siehe B7).
        \end{itemize}
\end{enumerate}

\subsubsection{Semistrukturiertes Interview}

Das Interview findet im Anschluss an das Post-Questionnaire statt und nutzt die Videoaufnahme des Playtests, um gezielt auf die während der Beobachtung notierten Auffälligkeiten einzugehen. Die folgenden Fragen dienen als Leitfaden und wurden je nach den tatsächlich beobachteten Auffälligkeiten der Testperson angepasst und ergänzt.

\begin{enumerate}[label=\textbf{I\arabic*}, leftmargin=1.5cm, itemsep=1em]
    \item
    \begin{enumerate}[label=\alph*), itemsep=0.6em]
        \item Du bist manchmal gegen die Wand gelaufen und hast versucht, dich daran zu greifen. Es hat öfters nicht funktioniert. Was hast du dir dabei überlegt? \\
            \textit{Antwort:} Ich hatte den Jump nicht ganz gecheckt, und ich wurde manchmal ungeduldig beim Springen und habe Ledge Grab versucht. Ich habe aber nicht ganz verstanden, dass ich gegen die Wand muss, und dass ich einfach Pech hatte, und habe es einfach noch einmal versucht. \\
            \textit{Bezug zur Beobachtung (B5):} Wiederholte Ledge-Grab-Versuche gegen die Wand statt gegen die greifbare Box; Sprünge gelangen anfangs schlecht und später etwas besser.
        \item Hast du die neuen Gegner immer nur anhand der Farbe erkannt, und konntest du an der Farbe den Typ erkennen? \\
            \textit{Antwort:} Ja, auch an den Attacken und an den Waffen, also z.\,B. an den verschiedenen Pfeilen. Und das Colorscheme, dass Archers so gelblich sind, war verständlich. \\
            \textit{Bezug zur Beobachtung (B9):} Beim ersten Heavy Archer erkannte sie sofort die zielsuchenden Geschosse; bei neuen Gegnern hielt sie meist kurz inne.
    \end{enumerate}

\item Du hattest zwei starke Reaktionen: einmal, als du Low Health hattest und viele Gegner da waren, und einmal, als du nicht hochgekommen bist. Was ist dir dabei durch den Kopf gegangen? \\
    \textit{Antwort:} Ich dachte eigentlich, dass ich mal sterben will, um zu sehen, was passiert, und dann habe ich durchgezogen. Beim Nicht-Hochkommen dachte ich mir \enquote{Cooked, scheisse, dass ich es nicht geschafft habe}, und dann habe ich mich resettet, es nochmal versucht und geschafft. \\
    \textit{Bezug zur Beobachtung (B2, B5):} Bei zu vielen Gegnern rannte sie weg, um zu heilen, und machte danach weiter; beim Hochkommen scheiterte sie wiederholt am Ledge Grab.

\item An dieser Stelle bist du wiederholt gestorben bzw. hast du gezögert. Was hat dir das Weiterkommen erschwert? \\
    \textit{Antwort:} Die Testperson ist nicht gestorben, die Frage entfällt.

\item Wie haben sich die grossen und kleinen Räume für dich unterschieden? \\
    \textit{Antwort:} Mein Lieblingsraum war Raum 7.5, weil er eine schöne neue Struktur eingeführt hat, nämlich die Brücke. Damit kann man mehr machen als mit nur zwei Containern. Es gab eine coole Abwechslung zwischen grossen und kleinen Räumen, oder besser gesagt eine Mischung. Kleine Räume waren eher schwerer, weil es keine Zuflucht gab. \\
    \textit{Bezug zur Beobachtung (B6, B12):} In kleinen Räumen spielte sie taktischer (Treffer austeilen und zurückziehen) und war genervt, wenn sie beim Heilen unterbrochen wurde.

\item Als hier mehrere Gegner gleichzeitig da waren: Wie hast du entschieden, wen du zuerst angreifst? \\
    \textit{Antwort:} Ich bin fast immer zuerst auf Archers gegangen, wenn ich im Nahkampf war. Vor allem die, die drei Geschosse oder suchende Pfeile geschossen haben, waren krass, deswegen wollte ich die immer zuerst loswerden. Die mit den Hämmern konnte man gut countern, und sie waren langsam, deswegen habe ich sie zum Schluss aufbewahrt. Und die Soldiers haben mich vom Heilen abgehalten, weil sie schnell waren, und deswegen habe ich sie nach den Archers getötet. \\
    \textit{Bezug zur Beobachtung (B3):} Sie griff meistens zuerst die Archers an, bei niedriger Gesundheit und nach Pfeiltreffern verstärkt; oft standen ihr dabei Soldiers im Weg.

\item Empfandest du diesen Treffer bzw. diesen Tod als eigenen Fehler oder als unfair? Warum? \\
    \textit{Antwort:} Wenn ich in der Luft getroffen wurde, war es deprimierend, bis ich wieder aufgestanden bin. War aber witzig. Homing Missiles waren schwerer zu dodgen, vor allem wenn man sie im letzten Moment bemerkt hat, und sie haben viel Schaden gemacht. Aber die Schüsse fühlten sich nicht fehl am Platz an und waren eher meine Fehler. Wenn viele Soldiers auf einmal gespawnt sind, war es ein wenig schwerer, lol, aber nicht unangenehm. \\
    \textit{Bezug zur Beobachtung (B4, B13):} Sie reagierte auf Angriffe meist erst nach dem Treffer und tankte viele Treffer; sie setzte manchmal unnötig Air Combat ein.

\item Gab es einen Punkt, an dem sich der Kampf für dich in einen anderen \enquote{Modus} verschoben hat (z.\,B. von entspannt zu angespannt)? Wo genau? \\
    \textit{Antwort:} Ich fand, dass Heilpausen schwer zu finden waren, und ich musste einfach wegrennen, um das machen zu können. Die Waves kamen mehr und mehr, und man musste sich einen Moment erarbeiten oder auf Kisten stehen. Es fühlte sich an wie überlappende Wellen. Wenn die Räume fertig waren, konnte man entspannen, aber in einem Raum gab es keine allzu guten Möglichkeiten. Ausser im Trap Room, dort wusste ich, dass es einer werden würde. \\
    \textit{Bezug zur Beobachtung (B7, B14):} Sie heilte an Orten ohne Gegner oder auf Boxen und erholte sich vor dem Raumwechsel; beim Raumwechsel bemerkte sie nicht immer gleich, dass sich die Tür geöffnet hatte.

\item Gibt es etwas, das du während des Spielens gedacht oder gefühlt hast, das wir bisher nicht angesprochen haben? \\
    \textit{Antwort:} Nichts, das mir jetzt einfallen würde. Also wir hatten kein Tutorial, und du hast es mir erklären müssen. Das wäre noch ein Verbesserungsvorschlag. \\
    \textit{Bezug zur Beobachtung (B10):} Sie bemerkte nach Abschluss eines Raums nicht gleich, dass sie weitergehen konnte.
\end{enumerate}

\subsection{Testperson 3 Bastian}

\subsubsection{Pre-Questionnaire}

Das Pre-Questionnaire dient dazu, die Testperson zu charakterisieren. Ihre Aussagen können so im Kontext gesehen und gewichtet werden.

\begin{enumerate}[label=\textbf{P\arabic*}, leftmargin=1.5cm, itemsep=0.8em]
\item Wie viele Stunden spielst du pro Woche Videospiele? \\
    \textit{Antwort:} 0 - 20h (kommt auf releases von Spielen an)

\item Was für Videospielgenres spielst du besonders gerne? \\
    \textit{Antwort:} Action, Platformer

\item Was macht für dich einen guten Kampfabschnitt aus? \\
    \textit{Antwort:} gutes Feedback; Dynamik; gut abgetrennte / erkennbare Segmente von Angriff und Verteidigung

\item Was stellst du dir unter einem Hack and Slash vor? \\
    \textit{Antwort:} Action hauptsächlich basierend auf Kampf; Spielzyklus mit ständig auftretenden neuen Erweiterungen; Steigerung der Fähigkeiten inner- und außerhalb des Spiels

\item Wie schätzt du dein eigenes Können in actionlastigen Spielen generell ein? (Skala 1--5: \enquote{Anfänger} bis \enquote{sehr erfahren}) \\
    \textit{Antwort:} 3

\item Bevorzugst du eher vorsichtiges, methodisches Spiel oder aggressives, schnelles Spiel? \\
    \textit{Antwort:} meist mit wenig Denken oder Konzipierung eines Spielstils, ich spiele wie ich mich fühle.

\item Was erwartest du dir von diesem Playtest, bevor du ihn spielst? \\
    \textit{Antwort:} demo-artiges Schülerprojekt (bereits ein Teil gesehen)

\item Worauf reagierst du in Kämpfen erfahrungsgemäss am empfindlichsten: z.\,B. unfaire Treffer, Unübersichtlichkeit, zu leichte Gegner, zu schwere Gegner? \\
    \textit{Antwort:} leichteres Design gefällt tendenziell besser; aber keine besondere Präferenz oder Abneigung
\end{enumerate}

\subsubsection{Post-Questionnaire}

Das Post-Questionnaire erfasst die subjektive Einschätzung der Testperson direkt im Anschluss an den Playtest, pro Designvorschlag sowie für den gesamten Kampfabschnitt. Hinter jeder Aussage steht in Klammern der Designvorschlag, den sie überprüft. Alle Aussagen werden auf einer Skala von 1 (\enquote{trifft gar nicht zu}) bis 5 (\enquote{trifft völlig zu}) bewertet, die offenen Fragen frei beantwortet.

\paragraph{Gegnerverhalten}
\begin{enumerate}[label=\textbf{Q\arabic*}, leftmargin=1.5cm, itemsep=0.8em]
\item Ich fühlte mich meinen Gegnern grundsätzlich überlegen, ausser in einzelnen zugespitzten Momenten. \hfill (V1) \\
    \textit{Antwort:} 5
\item Auch bei vielen Gegnern gleichzeitig behielt ich die Kontrolle über die Situation. \hfill (V2) \\
    \textit{Antwort:} 5
\item Ich musste bewusst entscheiden, welchen Gegner ich zuerst angreife. \hfill (V3) \\
    \textit{Antwort:} 5
\item Angriffe der Gegner wurden klar angekündigt, bevor sie mich trafen. \hfill (V4) \\
    \textit{Antwort:} 5
\end{enumerate}

\paragraph{Gegnerplatzierung und Kampfraum}
\begin{enumerate}[label=\textbf{Q\arabic*}, leftmargin=1.5cm, itemsep=0.8em, resume]
\item Die Gegner waren zu Beginn klar überschaubar positioniert, das Verhalten wurde danach unvorhersehbarer. \hfill (V5) \\
    \textit{Antwort:} 4
\item Der enge Abschnitt fühlte sich beim Spielen anders an als der offene Abschnitt. \hfill (V6) \\
    \textit{Antwort:} 3
\item Der Kampf hatte spürbare Phasen von Anspannung und Erholung. \hfill (V7) \\
    \textit{Antwort:} 5
\item Neue Gegnertypen wurden mir zuerst einzeln, dann erst in Kombination vorgestellt. \hfill (V8) \\
    \textit{Antwort:} 5
\end{enumerate}

\paragraph{Gesamteindruck}
\begin{enumerate}[label=\textbf{Q\arabic*}, leftmargin=1.5cm, itemsep=0.8em, resume]
\item Insgesamt war dieser Kampfabschnitt für mich befriedigend zu spielen. \hfill (Spielgefühl) \\
    \textit{Antwort:} 4
\item Der gesamte Kampfabschnitt fühlte sich zusammenhängend an, nicht wie eine Aneinanderreihung einzelner Momente. \hfill (Spielfluss) \\
    \textit{Antwort:} 1 für raum zu raum und 4 für locked in
\item Der Schwierigkeitsgrad war insgesamt angemessen. \hfill (Schwierigkeit) \\
    \textit{Antwort:} 5
\item Was hat sich am besten angefühlt? \hfill (offen) \\
    \textit{Antwort:} Das taktische wen angreifen, healen, weglaufen, homing missiles kan man cool ausweichen. Der schadensunterschied zwischen links und rechts click dass man den archer mit rechtsclick instantly killen kann
\item Was hat sich unfair oder unangenehm angefühlt? \hfill (offen) \\
    \textit{Antwort:} die bewegungen sind schwerfällig und auch wenn ich das telegraphing sehe kann ich nicht immer reagieren. die range konnte ich nicht von anfang an einschätzen und es hatte eine learning curve.
\item Wie war dieser Kampfabschnitt im Vergleich zu dem, was du dir vor dem Test erwartet hast? \hfill (offen) \\
    \textit{Antwort:} Die bewegung hätte ich ein wenig flüssiger erwartet, ansonsten hat es ziemlich viel mit dem vorwissen und von dem was ich mitbekommen habe übereingestimmt.
\end{enumerate}

\subsubsection{Beobachtungsraster}

Während des Playtests wird das Verhalten der Testperson anhand der folgenden Kriterien dokumentiert, jeweils mit Zeitstempel, um die Beobachtung später dem passenden Segment und dem retrospektiven Interview zuordnen zu können. Zunächst wird nur beobachtet und protokolliert, nicht interpretiert.

\paragraph{Gegnerverhalten}
\begin{enumerate}[label=\textbf{B\arabic*}, leftmargin=1.5cm, itemsep=0.8em]
\item Verhält sich die Testperson in normalen Phasen risikofreudig und aggressiv (z.\,B. wenig Ausweichen, viel Angreifen) und spürbar vorsichtiger in den zugespitzten Momenten? \hfill (V1)
    \begin{itemize}[nosep]
        \item Die Testperson ist ziemlich aggressiv, bewegt sich recht schnell und kommt zügig voran.
        \item Sie ignoriert ihr Low Health.
    \end{itemize}

\item Verhält sich die Testperson in Momenten mit vielen Gegnern anders (hektischer, mehr Fehler) als in ruhigeren Momenten? \hfill (V2)
    \begin{itemize}[nosep]
        \item Zum ersten Mal rennt sie vor einer Dreiergruppe Soldiers weg, um zu heilen.
        \item Sie gerät bei vielen Gegnern nicht in Panik, sondern nur, wenn sie wenig Leben hat und viele Gegner da sind.
    \end{itemize}

\item Welchen Gegnertyp greift die Testperson zuerst an, wenn Nah- und Fernkämpfer gleichzeitig anwesend sind? \hfill (V3)
    \begin{itemize}[nosep]
        \item Momentan geht sie einfach von Archer zu Archer, lässt die Soldiers sie verfolgen und schlägt diese nebenbei mit.
        \item Den normalen Soldier empfindet sie als schwerer bzw. nerviger als den Heavy Soldier, weil der Heavy Soldier weniger oft trifft. Die Heavy Soldiers haben deshalb die kleinste Priorität.
        \item Sie greift oft an, was sie gerade vor sich sieht (siehe B6).
    \end{itemize}

\item Reagiert die Testperson (z.\,B. durch Ausweichen), bevor ein Angriff trifft, oder erst danach bzw. gar nicht? \hfill (V4)
    \begin{itemize}[nosep]
        \item Sie versucht schon von Anfang an, Pfeilen auszuweichen, hält aber meistens einfach drauf und ignoriert den Schaden.
        \item Sie reagiert momentan selten auf Angriffe, besonders auf die der Nahkämpfer. Erst wenn nur noch Archers übrig sind, beginnt sie auszuweichen.
        \item Sie will zunehmend mehr ausweichen, kann es aber nicht gut, weil ihre eigenen Combos zu schwerfällig sind und ihr die Zeit zum Ausweichen fehlt.
        \item Sie ist generell gut in der Luftbewegung (Aerial Movement) und weicht dort oft aus.
    \end{itemize}

\item Wiederholt die Testperson denselben Fehler mehrfach an derselben Stelle?
    \begin{itemize}[nosep]
        \item Sie wird wiederholt in der Luft von Archers abgeschossen und fällt zu Boden.
        \item Sie springt weiterhin während des Kampfes, wahrscheinlich aus Gewohnheit (Muscle Memory) aus anderen Spielen. In diesem Spiel bringt das nichts.
        \item Das Ledge Grab beherrscht sie solide. Sie muss nicht pixelgenau springen, um auf die Box zu kommen, weil sie es kann und es weiss.
        \item Zur Luftkombo siehe B13.
    \end{itemize}
\end{enumerate}

\paragraph{Gegnerplatzierung und Kampfraum}
\begin{enumerate}[label=\textbf{B\arabic*}, leftmargin=1.5cm, itemsep=0.8em, resume]
\item Zeigt die Testperson zu Beginn eines Encounters ein zielgerichtetes Vorgehen (z.\,B. gezieltes Anvisieren einzelner Gegner) und wird dieses Vorgehen unruhiger, sobald sich die Gegner umzingelnd verhalten? \hfill (V5)
    \begin{itemize}[nosep]
        \item Sie greift oft den Gegner an, den sie gerade vor sich sieht.
        \item Einmal wird sie in der Luft von einem Archer getroffen, fällt zu Boden und wird dort von Soldiers und Heavy Soldiers umzingelt. Sie wird dafür deutlich bestraft.
        \item Sie findet einen versteckten Heavy Archer nicht und sucht ihn. Der Heavy Archer hat sie auch nicht aus dem Hinterhalt getroffen, deshalb hat sie ihn nicht bemerkt.
    \end{itemize}

\item Unterscheidet sich das Bewegungsmuster im engen Abschnitt sichtbar vom offenen Abschnitt (mehr Ausweichen, mehr Distanzhalten)? \hfill (V6)
    \begin{itemize}[nosep]
        \item In grossen sowie in kleinen Räumen ist sie genau gleich aggressiv.
        \item Zu einem anderen Zeitpunkt ist sie in kleinen Räumen hektischer als in grossen und versucht, die Gegner schneller zu töten.
    \end{itemize}

\item Zeigt die Testperson in Erholungsphasen ein erkennbar anderes Verhalten (z.\,B. Innehalten, freies Erkunden) als während einer Welle? \hfill (V7)
    \begin{itemize}[nosep]
        \item Sie stirbt einmal, weil sie nie geheilt hat.
        \item Zum ersten Mal heilt sie.
        \item Zwischen den Räumen macht sie kaum Pausen zum Heilen und geht einfach weiter.
        \item In Ruhephasen heilt sie inzwischen regelmässiger.
        \item Sie nimmt gern Abstand, um zu heilen. Das hat sie inzwischen gelernt (Eindruck der Beobachtenden).
        \item Sie hält eine ruhige Phase für das Ende des Raums (siehe B11).
        \item Zur Fluchtreaktion statt Box siehe B13.
    \end{itemize}

\item Wirkt die Testperson beim ersten Auftreten eines neuen Gegnertyps zögerlich oder unsicher, und verschwindet dieses Zögern bei späteren Wiederholungen desselben Typs? \hfill (V8)
    \begin{itemize}[nosep]
        \item Der Heavy Soldier interessiert sie nicht, sie greift ihn einfach an. Sie hat vorher gesagt, dass sie lieber ohne viel Nachdenken spielt (vgl. P6).
        \item Beim ersten Heavy Archer mit Homing Missile versucht sie auszuweichen, schafft es nicht und tankt den Treffer. Es scheint ihr egal zu sein.
        \item Den Chain Archer sieht sie zum ersten Mal. Es interessiert sie wenig, sie erkennt aber, dass es eine neue Art Gegner ist.
    \end{itemize}
\end{enumerate}

\paragraph{Allgemein}
\begin{enumerate}[label=\textbf{B\arabic*}, leftmargin=1.5cm, itemsep=0.8em, resume]
\item Anzeichen von Verwirrung oder Orientierungslosigkeit (z.\,B. Stehenbleiben, Umherblicken ohne erkennbares Ziel)
    \begin{itemize}[nosep]
        \item Etwa in der Mitte des Spiels benutzt sie zum ersten Mal richtig den Double Jump. Es verwirrt sie, dass sie lange drücken kann, um höher zu springen, und sie nutzt es deshalb nicht. Sie spammt die Leertaste, möglicherweise weil sie etwas gestresst ist (Vermutung der Beobachtenden).
        \item Sie will in den nächsten Raum, kann aber nicht, und sucht deshalb den restlichen Gegner.
        \item Sie sucht den versteckten Heavy Archer (siehe B6).
    \end{itemize}

\item Spontane verbale Reaktionen (Wortlaut und Zeitstempel notieren)
    \begin{itemize}[nosep]
        \item Sie sagt, sie glaube, es sei die letzte Welle des Raums, obwohl sie wenig Leben hat. Es folgt aber noch eine weitere Welle. Weil es so ruhig war, dachte sie, es sei die letzte gewesen, hat sich aber getäuscht.
        \item Im zweitletzten Raum sagt sie \enquote{Boss} (laut Interview, siehe I2).
    \end{itemize}

\item Sichtbare Frustration oder sichtbare Freude/Erleichterung, mit Zeitstempel
    \begin{itemize}[nosep]
        \item Sie wirkt sehr fokussiert.
        \item Sie hat sichtlich Spass und lacht die ganze Zeit.
        \item Besonders der Heavy Archer nervt sie.
        \item Einmal wird sie von Gegnern umzingelt und getroffen und ist sichtlich genervt (siehe I2).
        \item Nach dem Ende des Spiels spielt sie es direkt von vorn.
    \end{itemize}

\item Auffällige, nicht vorhergesehene Entscheidungen (z.\,B. Ignorieren eines Gegnertyps, unerwartete Fluchtversuche)
    \begin{itemize}[nosep]
        \item Sie versucht oft die Luftkombo, obwohl sie damit die Bodengegner seltener trifft.
        \item Statt auf eine Box zu steigen und Pause zu machen, versucht sie wegzurennen. Die Beobachtenden verstehen das nicht ganz, vielleicht hat sie es vergessen.
        \item Ihre momentane Taktik: Gegner anhäufen, hineinspringen, etwas Schaden machen und wieder herausgehen, immer wieder.
    \end{itemize}

\item Wirkt der Übergang zwischen zwei Segmenten (z.\,B. Welle zu Erholung, eng zu offen) beim Spielverhalten fliessend, oder gibt es einen sichtbaren Bruch (Zögern, Pause, sichtbare Neuorientierung)? \hfill (Spielfluss)
    \begin{itemize}[nosep]
        \item Sie macht zwischen den Räumen kaum Pausen und geht einfach weiter (siehe B8).
        \item Sie will in den nächsten Raum, kann aber nicht, und sucht deshalb den restlichen Gegner (siehe B10).
    \end{itemize}
\end{enumerate}

\subsubsection{Semistrukturiertes Interview}

Das Interview findet im Anschluss an das Post-Questionnaire statt und nutzt die Videoaufnahme des Playtests, um gezielt auf die während der Beobachtung notierten Auffälligkeiten einzugehen. Die folgenden Fragen dienen als Leitfaden und wurden je nach den tatsächlich beobachteten Auffälligkeiten der Testperson angepasst und ergänzt.

\begin{enumerate}[label=\textbf{I\arabic*}, leftmargin=1.5cm, itemsep=1em]
\item
\begin{enumerate}[label=\alph*), itemsep=0.6em]
    \item Was hast du gedacht, bevor du zum ersten Mal gestorben bist? \\
        \textit{Antwort:} Ich hatte Heal vergessen und Leben ignoriert. \\
        \textit{Bezug zur Beobachtung (B1, B8):} Sie ignorierte ihr Low Health und starb, weil sie nie geheilt hatte.
    \item Wie hast du die neuen Gegner wahrgenommen? \\
        \textit{Antwort:} Beim ersten Archer dachte ich, es ist einfach eine stärkere Version, und die anderen Archers habe ich auch nur am Angriff erkannt. Sonst habe ich die anderen schon bemerkt, wusste aber nicht genau, was sie machen. \\
        \textit{Bezug zur Beobachtung (B9):} Beim Chain Archer zeigte sie wenig Interesse, erkannte aber, dass es eine neue Art Gegner ist.
\end{enumerate}

\item
\begin{enumerate}[label=\alph*), itemsep=0.6em]
\item Warum hast du im zweitletzten Raum \enquote{Boss} gesagt? \\
    \textit{Antwort:} Ja halt, grosser Raum und leer = Boss. \\
    \textit{Bezug zur Beobachtung (B11):} Verbale Reaktion im zweitletzten Raum.
\item Du warst einmal von Gegnern umgeben, wurdest getroffen und warst sichtlich genervt. Warum? \\
    \textit{Antwort:} Ich war stuck, es gab gefühlt einen Input Delay, und ich habe auf die Gegner nicht geachtet, und es hat mich genervt. \\
    \textit{Bezug zur Beobachtung (B4, B12):} Ihre eigenen Combos waren zu schwerfällig, um rechtzeitig auszuweichen; sichtbare Frustration beim Umzingeltwerden.
\end{enumerate}

\item An dieser Stelle bist du wiederholt gestorben bzw. hast du gezögert. Was hat dir das Weiterkommen erschwert? \\
\textit{Antwort:} Die Testperson ist nicht wiederholt gestorben, die Frage entfällt.

\item Wie unterscheiden sich die kleinen, mittleren und grossen Räume für dich? \\
\textit{Antwort:} Bei den kleinen Räumen wusste ich direkt, wohin. Bei den Räumen mit den Boxen hatte ich besonders viele Optionen, und der Raum mit der Brücke hatte am meisten Optionen. Gestresst war ich nicht, weil ich kontrollieren konnte, ob ich rein- oder rausgehe. \\
\textit{Bezug zur Beobachtung (B7):} Sie war in grossen und kleinen Räumen gleich aggressiv, zu einem anderen Zeitpunkt in kleinen Räumen hektischer.

\item Als hier mehrere Gegner gleichzeitig da waren: Wie hast du entschieden, wen du zuerst angreifst? \\
\textit{Antwort:} Definitiv erst mal die Archers, vor allem die normalen, danach die anderen Archer-Typen. Die waren halt immer omnipräsent, und ich habe viel Schaden kassiert. Die Nahkampfgegner kommen sowieso zu einem, und es gibt keinen Unterschied zwischen gross und klein. Die kleineren hatten nicht wirklich Präsenz. Die grossen hatten Momente, in denen man dachte, ich muss ausweichen, aber die haben auch oft nicht getroffen. \\
\textit{Bezug zur Beobachtung (B3):} Sie ging von Archer zu Archer und liess die Soldiers folgen; Heavy Soldiers hatten die kleinste Priorität.

\item Empfandest du diesen Treffer bzw. diesen Tod als eigenen Fehler oder als unfair? Warum? \\
\textit{Antwort:} Ja, nach den Attack-Combos sieht man zum Beispiel, wie Arrows kommen, und man kann trotzdem nicht ausweichen, weil der Input nicht freigegeben wurde. Aber man hatte genug Healing-Möglichkeiten. \\
\textit{Bezug zur Beobachtung (B4, B5):} Ihre eigenen Combos waren zu schwerfällig, um rechtzeitig auszuweichen; sie wurde wiederholt in der Luft abgeschossen.

\item Gab es einen Punkt, an dem sich der Kampf für dich in einen anderen \enquote{Modus} verschoben hat (z.\,B. von entspannt zu angespannt)? Wo genau? \\
\textit{Antwort:} Nachdem man alle Ranged Enemies getötet hat, war es immer entspannt. Davor ist es recht hektisch, und ich muss aufpassen. Aber es ergab eine gute Balance von Spannung und Entspannung. Am Anfang Spannung, am Ende Entspannung. Bugs zu finden hat mir besonders Freude bereitet. \\
\textit{Bezug zur Beobachtung (B8, B11, B14):} In Ruhephasen heilte sie inzwischen regelmässiger; sie hielt eine ruhige Phase fälschlich für das Raumende; zwischen den Räumen machte sie kaum Pausen.

\item Wie hast du das Erscheinen der Gegner und ihre Platzierung wahrgenommen? \\
\textit{Antwort:} Die Gegner sind zwar ausserhalb meiner Sichtweite gespawnt, man hat sie aber noch früh genug bemerkt. Ich glaube, die Gegnerplatzierung ist stellenweise festgesetzt, und danach spawnen halt Archers eher erhöht und in Gruppen. Also die erste Welle ist definitiv immer gesetzt, zumindest teilweise. \\
\textit{Bezug zur Beobachtung (B6):} Sie griff meist den Gegner vor sich an; einen versteckten Heavy Archer fand sie nicht.
\end{enumerate}

\subsection{Testperson 4 Max}
\subsubsection{Pre-Questionnaire}

Das Pre-Questionnaire dient dazu, die Testperson zu charakterisieren. Ihre Aussagen können so im Kontext gesehen und gewichtet werden.

\begin{enumerate}[label=\textbf{P\arabic*}, leftmargin=1.5cm, itemsep=0.5em]
\item Früher mehrere Stunden, heutzutage ist es ziemlich genau 0h. :(
\item Am liebsten sind Taktik-Spiele oder Zivilisationsaufbauspiele, mit Kriegselementen. Jedoch ist auch ab und zu ein Spiel, welches sich nur auf Geballer konzentriert, nett.
\item Gute Visualisierung (Sprich, Gegner muss sichtbar sein, das Umfeld soll schön sein und zur atmosphäre beitragen…), ein smoothes Handling und befriedigende Combos.
\item Unter einem HAck and Slash Spiel, stelle ich mir etwas vor, dass nicht viel Hirnleistung erfordert und einfach da ist um cool auszusehen, bzw. dich gut fühlen zu lassen mit spektakulären combos etc. pp.
\item 3
\item Beides hat seine Reize, wenn ich müde bin, dann zweiteres, sonst möchte ich Hirnfutter.
\item Ein mediocre af game, was aber spass macht und an manchen stellen unnötig krass gut sein wird, an anderen wiederum den anschein erweckt, 3 Tage lang gemacht worden zu sein.
\item Unsichtbare Feinde ohne jegliche Andeutung machen keinen Sinn. Und einfach 3 Minuten einem Gegner ausweichen zu müssen, ohne eine Gegenspiel Möglichkeit zu haben ist auch wack. Zu schwierige Gegener sind probably auch nervig.
\end{enumerate}

\subsubsection{Post-Questionnaire}

Das Post-Questionnaire erfasst die subjektive Einschätzung der Testperson direkt im Anschluss an den Playtest, pro Designvorschlag sowie für den gesamten Kampfabschnitt. Hinter jeder Aussage steht in Klammern der Designvorschlag, den sie überprüft. Alle Aussagen werden auf einer Skala von 1 (\enquote{trifft gar nicht zu}) bis 5 (\enquote{trifft völlig zu}) bewertet, die offenen Fragen frei beantwortet.

\paragraph{Gegnerverhalten}
\begin{enumerate}[label=\textbf{Q\arabic*}, leftmargin=1.5cm, itemsep=0.5em]
\item 5
\item 4
\item 3
\item 1
\end{enumerate}

\paragraph{Gegnerplatzierung und Kampfraum}
\begin{enumerate}[label=\textbf{Q\arabic*}, leftmargin=1.5cm, itemsep=0.5em, resume]
\item 2
\item 2
\item 3
\item 5
\end{enumerate}

\paragraph{Gesamteindruck}
\begin{enumerate}[label=\textbf{Q\arabic*}, leftmargin=1.5cm, itemsep=0.5em, resume]
\item 4
\item 4
\item 5
\item Horden an kleinen Viechern zu zerstören. Das movement war sehr cool gemacht, hatte aber keine wirklichen Auswirkungen
\item Die Combos haben sehr gestört, da nie ganz klar war wann ich wieder angreifen kann oder verwundbar bin, zudem konnten sie nicht abgebrochen werden.
\item Ähnlich zum erwartetem, wobei bei der Variation der Gegner positiv überrascht, blocken und weitere Kontrollierbare Elemente hätte ich aber schon erwartet.
\end{enumerate}

\subsubsection{Beobachtungsraster}
Max benutzt die Flying combo ohne grund.
max ignoriert die gegner und probiert das movement aus in raum 1.
max hat nicht gecheckt dass die türe offen ist.
max hat ohne nachzudenken auf den neu eingeführten archer eingeschlagen.
max weicht den projektilen aus und hält abstand zu den archern
max benutzt vor allem die light attack combo
max versucht mit objekten zu interagieren.
max war erstaunt als er in den 3. raum gekommen ist.
max spammt dashen statt es gedrückt zu halten.
max schlägt zuerst die gegner die auf ihn zukommen
max hat zum ersten mal gerade die heavy attack benutzt
max meinte, dass die gegner irgendwie nicht auf hits reagieren aber als er genau hingesehen hat, ist es ihm aufgefallen, als er genauer hingesehen hat.
max meinte es wäre cool wenn es einen so jump drop attack gäbe
max fuchtelt mit der kamera
max häuft soldiers an bevor er sie angreift, und er fängt an anzugreifen, wenn er bei einem archer steht um alles auf einmal schlagen zu können.
sobalt max eine wave sieht schaut er sie erst and bevor er angreift.
max ist irgendwie die ganze zeit auf halbem leben
max hat wiede rdie offene tür nicht bemerkt.
max hat bei einem angriff nicht bemerkt, dass er einen gegner geschlagen hat.
max wurde beim aufstehen nach einem schlag wieder geschlagen und konnte nichts machen.
max reagiert nicht auf schaden sondern läuft einfach rein und schlägt
max healt bevor er in den nächsten raum geht.
max dashed nicht viel
max springt gerne in so mengen rein, weil er meinte es gäbe eine sprungattacke.
max meint er geniesse seine überlegenheit.
max sagt, es macht eigentlich schon noch ganz spass.
max hat instantly gemeint, dass es ein neuer gegner ist beim heavy archer.
max meint es ist besonders befiredigend, wenn man alle auf einmal killen kann.
max meint, es ist etwas clunky mit dem angreifen aber es sei nicht so schilmm für unseren combat.
max wurde von einem heavy archers getroffen und konzentriert jetzt auf heavy archers.
max schlägt immernoch weiter auf den ersten
max meint das einzige das punished im game seien die etwas längeren kampfanimationen und wenn er in der luft getroffen wird.
max gemerkt, dass der chain archer neu ist und hat ihn nicht angegriffen um zu sehen was er macht.
max benutzt lieber die light attack weil er meint die heavy attack mache zu wenig damage für wie lange sie geht.
max findet besonders den heavy archer schlimm, weil er von dem getroffen wird aber ihn jucken die chain archers nicht.
max meint er würde lieber die combos abbrechen können aber er meint er checke wieso das so designed ist.
max ist gestorben in raum 9 weil er den archer nicht gesehen hat und getroffen wurde
nachdem max gestorben ist, macht er schneller die räume durch
also er spielt jetzt deutlich aggressiver.
max findet es nicht cool wenn er von hinten angegriffen wird.
mittlerweile geht max zuerst auf die archers.
max findet es spassig, die gegner zu predicten, weil er ja die räume schon gesehen hat.
max findet die ästhtik schön.
max spielt weniger taktisch als mehr einfach anzugreifen.
max wünscht sich einen block.
max findet das movement richtig cool und findet, dass es eine lernkurve hat.
max versucht wenn es geht auszuweigen aber meistens läuft er einfach durch.
max kann nicht genau einschätzen wie lange das heilen dauert.
max springt in der luft und versucht sich in der luft zu positionieren bevor er in die gegnerwelle eintaucht.
max geht eher auf die soldiers als auf die heavy solderis weil die wahrscheinlich dass diese ihn treffen kleinder ist.
max schaut eigentlich gar nicht mehr beim schalgen, weil er egal was passiert tifft.
max findet die animationen vom heavy soldier und chain archer cool.

Während des Playtests wird das Verhalten der Testperson anhand der folgenden Kriterien dokumentiert, jeweils mit Zeitstempel, um die Beobachtung später dem passenden Segment und dem retrospektiven Interview zuordnen zu können. Zunächst wird nur beobachtet und protokolliert, nicht interpretiert.

\paragraph{Gegnerverhalten}
\begin{enumerate}[label=\textbf{B\arabic*}, leftmargin=1.5cm, itemsep=0.5em]
\item Verhält sich die Testperson in normalen Phasen risikofreudig und aggressiv (z.\,B. wenig Ausweichen, viel Angreifen) und spürbar vorsichtiger in den zugespitzten Momenten? \hfill (V1)
\item Verhält sich die Testperson in Momenten mit vielen Gegnern anders (hektischer, mehr Fehler) als in ruhigeren Momenten? \hfill (V2)
\item Welchen Gegnertyp greift die Testperson zuerst an, wenn Nah- und Fernkämpfer gleichzeitig anwesend sind? \hfill (V3)
\item Reagiert die Testperson (z.\,B. durch Ausweichen), bevor ein Angriff trifft, oder erst danach bzw. gar nicht? \hfill (V4)
\item Wiederholt die Testperson denselben Fehler mehrfach an derselben Stelle?
\end{enumerate}

\paragraph{Gegnerplatzierung und Kampfraum}
\begin{enumerate}[label=\textbf{B\arabic*}, leftmargin=1.5cm, itemsep=0.5em, resume]
\item Zeigt die Testperson zu Beginn eines Encounters ein zielgerichtetes Vorgehen (z.\,B. gezieltes Anvisieren einzelner Gegner) und wird dieses Vorgehen unruhiger, sobald sich die Gegner umzingelnd verhalten? \hfill (V5)
\item Unterscheidet sich das Bewegungsmuster im engen Abschnitt sichtbar vom offenen Abschnitt (mehr Ausweichen, mehr Distanzhalten)? \hfill (V6)
\item Zeigt die Testperson in Erholungsphasen ein erkennbar anderes Verhalten (z.\,B. Innehalten, freies Erkunden) als während einer Welle? \hfill (V7)
\item Wirkt die Testperson beim ersten Auftreten eines neuen Gegnertyps zögerlich oder unsicher, und verschwindet dieses Zögern bei späteren Wiederholungen desselben Typs? \hfill (V8)
\end{enumerate}

\paragraph{Allgemein}
\begin{enumerate}[label=\textbf{B\arabic*}, leftmargin=1.5cm, itemsep=0.5em, resume]
\item Anzeichen von Verwirrung oder Orientierungslosigkeit (z.\,B. Stehenbleiben, Umherblicken ohne erkennbares Ziel)
\item Spontane verbale Reaktionen (Wortlaut und Zeitstempel notieren)
\item Sichtbare Frustration oder sichtbare Freude/Erleichterung, mit Zeitstempel
\item Auffällige, nicht vorhergesehene Entscheidungen (z.\,B. Ignorieren eines Gegnertyps, unerwartete Fluchtversuche)
\item Wirkt der Übergang zwischen zwei Segmenten (z.\,B. Welle zu Erholung, eng zu offen) beim Spielverhalten fliessend, oder gibt es einen sichtbaren Bruch (Zögern, Pause, sichtbare Neuorientierung)? \hfill (Spielfluss)
\end{enumerate}

\subsubsection{Semistrukturiertes Interview}

Das Interview findet im Anschluss an das Post-Questionnaire statt und nutzt die Videoaufnahme des Playtests, um gezielt auf die während der Beobachtung notierten Auffälligkeiten einzugehen. Die folgenden Fragen dienen als Leitfaden und werden je nach den tatsächlich beobachteten Auffälligkeiten der jeweiligen Testperson angepasst und ergänzt.

\begin{enumerate}[label=\textbf{I\arabic*}, leftmargin=1.5cm, itemsep=0.5em]
\item Als du zum zweiten Mal gestorben bist: Ja war halt dumm gelaufen mti den Combos, welche einfach ned aufgehört haben.

Was ist dir bei Air Combos durchgegangen?: Useless af, es gab halt dann einfach nur nervige Situationen wo ich in der Luft gehangen bin.

\item Heavysoldier?: Ja ist ein bisschen ein sponge was schaden angeht, war aber mehr lustig dass er die Hälfte der Zeit verfehlt hat und sonst nur in Gruppen zur Gefahr wurde was bissl konträr ist.
\item Beim drop in?: Ja ned ganz gecheckt, dass ich eigentlich nicht gepunished werde für grosse Höhen. Ansonsten war halt grosse Arena, aber nichts los ausser einem Lörach.
\item Zwei Zwsichenräume gr. kl.:? Für mich hat es keinen Unterschied gemacht, vllt war ich aber auch einfach nur im Ban von Hauen und Smashen.
\item Als hier mehrere Gegner gleichzeitig da waren: Ja wenn ich wollte Archer, weil die halt mies nerven und alle unterschiedlich angreifen ohne wirkliche prediktion oder sound effects die mich warnen könnten. Sonst auf jolo.
\item Empfandest du diesen Treffer bzw. diesen Tod als eigenen Fehler oder als unfair? Warum? Naja irgendwo skill issue, war aber auch blöd, da die Combos manchmal mies geback fired haben und ich nicht blocken konnte oder so shyyt.
\item Gab es einen Punkt, an dem sich der Kampf für dich in einen anderen \enquote{Modus} verschoben hat (z.\,B. von entspannt zu angespannt)? Wo genau?
Ich war eigentlich immer locked in und habe mich nur darauf konzentriert möglichst schnell die maus zu drücken, bin ich ehrlich herrr röörich.
\item Ja meine Geheimen Bestechungsgelder wurden nicht empfangen. NA FR ON GOD JETZT ich habe das Gefühl, dass man mehr Interaktionen hätte bauen sollen zum spielen, wobei button mashing schon funny ist, aber während dessen auch noch block oder dodge roll noch funnier wäre. YFM?
\end{enumerate}

\subsection{Testperson 5 Luca}
\subsubsection{Pre-Questionnaire}

Das Pre-Questionnaire dient dazu, die Testperson zu charakterisieren. Ihre Aussagen können so im Kontext gesehen und gewichtet werden.

\begin{enumerate}[label=\textbf{P\arabic*}, leftmargin=1.5cm, itemsep=0.5em]
\item je nach dem 4-30h
\item Singerlplayer open World, Factory games, multiplayer sandbox games, soulslike, rouglikes
\item challenging, gutes design, eigenes( gute controlles sei es movement, aim, parry), gut passende musikalische untermalung und soundeffekts
\item sehr viele gegner töten
\item In action games 4
\item methodisch, alles mögliche ausprobieren
\item Ich erwarte ein fertige fighting game mit einstimmer grafik gutem gamedesigne und das es spass macht, Gute setting für anpassungen
\item unfaire treffer/gegner, allgemein dinge die man nicht vorhersehen kann, dann aber dafür bestraft wird.
\end{enumerate}

\subsubsection{Post-Questionnaire}

Das Post-Questionnaire erfasst die subjektive Einschätzung der Testperson direkt im Anschluss an den Playtest, pro Designvorschlag sowie für den gesamten Kampfabschnitt. Hinter jeder Aussage steht in Klammern der Designvorschlag, den sie überprüft. Alle Aussagen werden auf einer Skala von 1 (\enquote{trifft gar nicht zu}) bis 5 (\enquote{trifft völlig zu}) bewertet, die offenen Fragen frei beantwortet.

\paragraph{Gegnerverhalten}
\begin{enumerate}[label=\textbf{Q\arabic*}, leftmargin=1.5cm, itemsep=0.5em]
\item 3
\item 4
\item 4
\item 3
\end{enumerate}

\paragraph{Gegnerplatzierung und Kampfraum}
\begin{enumerate}[label=\textbf{Q\arabic*}, leftmargin=1.5cm, itemsep=0.5em, resume]
\item 2
\item 5
\item 5
\item 5
\end{enumerate}

\paragraph{Gesamteindruck}
\begin{enumerate}[label=\textbf{Q\arabic*}, leftmargin=1.5cm, itemsep=0.5em, resume]
\item 4
\item 4
\item 5
\item Wenn du jemanden Spawnkillst, also wenn du schon weisst wo er spawnt oder du vermutest wo dieser spawn.
\item Wenn ich stunlocked worden bin von ganz vielen gegnern.
\item ich dachte ich wäere noch viel stärker also wie hyrule warriors und ich habe erwartet dass ich 1000de umbringe
\end{enumerate}

\subsubsection{Beobachtungsraster}
bei neuen gegnern greft luca an ohne zu schauen was sie machen.
luca versucht glitches zu finden.
luca ist erst einmal die klippe runter gesprungen um zu sehen was passiert.
er versucht durch die map zu klippen.
luca kann nicht ganz einschätzen was seine range ist.
er greift das an, was vor ihm ist.
luca will emoten können.
luca ignoriert die soliders und geht einfach auf die archers.
luca wollte die gegner groupen und dann killen aber hat dabei viel damage kassiert
luca steht auf die props um vor den gegnern zuflucht zu finden, das haben wir noch nie bei einem spieleg gesehen.
luca rennt im kreiss um sich taktiken zu überlegen. und während im kreiss rennen trifft ihn halt nix.
luca hat gecheckt ob gegner arrows abfangen können.
luca ist etwas anders als die anderen playtester, da er mehr taktisch als agressiv angeht.
luca versucht die gegner-ai zu checken.
luca wirkt sehr selbstsicher.
luca sucht auch nach interactables. es gibt aber nichts
luca sammelt die gegner bei einem archer und versucht alles zu killen aber schaft es nicht immer ganz so clean
luca findet es spielt sich wie hyrule warriors, was er gespielt hat und es ist das gleiche genre.
luca alternated zwischen heavy und light attack
luca findet archers nervig, weil sie ihn treffen aber er geht immernoch zuerst auf close range enemies.
luca heilt wenn nur noch wenige eneimes überbleiben am ende einer welle.
luca kämpft vor allem mit soldiers und greift nur hin und wieder einen archer an, wenn dieser ihn abgeschossen hat.
Luca greift auch den heavy soldier an ohne zu zögern
luca sucht den highground zum healen.
luca versucht die gegner aus seinem highground zu cheasen aber das geht nicht.
luca ist gestorben weil er zu wenig leben hatte und nicht mehr heilen konnte weil zuwenig mana
luca geht jetzt ziemlich schnell durch die räume durch aber immernoch eher taktisch.
luca findet die telegraphing animation der archers ein wenig verwirrent weil er meint es sieht aus als laufe er aber dieser bliebt stehen.
luca möchte so heavy und light attack knöpfe kombinieren zu combos aber das haben wir nicht wir haben nur halt kombos mit pause und so
luca merkt das es unterschiedliche schadenswerte macht und hat gemerkt das 2x heavy attack den archer dierekt killt.
luca greift zuerst das an, was auf den boxen steht um sich eine zuflucht zu sichern, falls es knapp werden sollte.
luca findet auch das der Heavysoldier einfach irgendwo hin springt.
luca hat ein wenig probleme mit dem movement also das switchen zwischen movement und combat.
wenn luca damage nimmt, zieht er sich zurück, andere playtester hätten einfach weitere drauf gehalten. das punisht ihn aber mehr weil, wenn sie alle einmal angegriffen haben, kann man sie dann easy killen. er aber wird angegriffen und verzieht sich deshalb kann er dann kein damage dealen.
luca ist daher eher sensibel auf schaden.
luca meinst er kann auch nicht einschätzen wie viel schaden so eine gruppe von gegnern macht, wenn er reinläuft.
luca meint er erkennt die verschiedenen und neuen gegner schon anhand der Frabe und des Aussehens.
luca ist die ganze zeit eher beschäftigt, das game zu beobechten, patterns zu finden und auch bugs.
luca meint, er wäre nie unter stress oder so.
luca meint der hammer sound pase irgendwie nicht lol
luca versucht jetzt gegner abzusondern und sie so danach anzugreifen weil er gruppenkämpfe nicht so einfach findet,also sie machen ihm zu viel schaden. aber das absondern ist nicht so einfach.
luca ist in raum 10 vorsichtig beim um die ecken gehen
luca hat keine probleme damit zu checken, dass die türen offen sind.
luca ist dierekt über dem drop in gewesen in raum 11 und ist dierekt in raum 12 gekommen. er war kurz gejumped aber hat dierekt den raum gecheckt und deshalb hat er sich dierekt beruhigt.
luca wäre gerade fast in einen türrahmen gedrückt worden. und wäre fast nicht mehr rausgekommen wegen der gegnermengen.

Während des Playtests wird das Verhalten der Testperson anhand der folgenden Kriterien dokumentiert, jeweils mit Zeitstempel, um die Beobachtung später dem passenden Segment und dem retrospektiven Interview zuordnen zu können. Zunächst wird nur beobachtet und protokolliert, nicht interpretiert.

\paragraph{Gegnerverhalten}
\begin{enumerate}[label=\textbf{B\arabic*}, leftmargin=1.5cm, itemsep=0.5em]
\item Verhält sich die Testperson in normalen Phasen risikofreudig und aggressiv (z.\,B. wenig Ausweichen, viel Angreifen) und spürbar vorsichtiger in den zugespitzten Momenten? \hfill (V1)
\item Verhält sich die Testperson in Momenten mit vielen Gegnern anders (hektischer, mehr Fehler) als in ruhigeren Momenten? \hfill (V2)
\item Welchen Gegnertyp greift die Testperson zuerst an, wenn Nah- und Fernkämpfer gleichzeitig anwesend sind? \hfill (V3)
\item Reagiert die Testperson (z.\,B. durch Ausweichen), bevor ein Angriff trifft, oder erst danach bzw. gar nicht? \hfill (V4)
\item Wiederholt die Testperson denselben Fehler mehrfach an derselben Stelle?
\end{enumerate}

\paragraph{Gegnerplatzierung und Kampfraum}
\begin{enumerate}[label=\textbf{B\arabic*}, leftmargin=1.5cm, itemsep=0.5em, resume]
\item Zeigt die Testperson zu Beginn eines Encounters ein zielgerichtetes Vorgehen (z.\,B. gezieltes Anvisieren einzelner Gegner) und wird dieses Vorgehen unruhiger, sobald sich die Gegner umzingelnd verhalten? \hfill (V5)
\item Unterscheidet sich das Bewegungsmuster im engen Abschnitt sichtbar vom offenen Abschnitt (mehr Ausweichen, mehr Distanzhalten)? \hfill (V6)
\item Zeigt die Testperson in Erholungsphasen ein erkennbar anderes Verhalten (z.\,B. Innehalten, freies Erkunden) als während einer Welle? \hfill (V7)
\item Wirkt die Testperson beim ersten Auftreten eines neuen Gegnertyps zögerlich oder unsicher, und verschwindet dieses Zögern bei späteren Wiederholungen desselben Typs? \hfill (V8)
\end{enumerate}

\paragraph{Allgemein}
\begin{enumerate}[label=\textbf{B\arabic*}, leftmargin=1.5cm, itemsep=0.5em, resume]
\item Anzeichen von Verwirrung oder Orientierungslosigkeit (z.\,B. Stehenbleiben, Umherblicken ohne erkennbares Ziel)
\item Spontane verbale Reaktionen (Wortlaut und Zeitstempel notieren)
\item Sichtbare Frustration oder sichtbare Freude/Erleichterung, mit Zeitstempel
\item Auffällige, nicht vorhergesehene Entscheidungen (z.\,B. Ignorieren eines Gegnertyps, unerwartete Fluchtversuche)
\item Wirkt der Übergang zwischen zwei Segmenten (z.\,B. Welle zu Erholung, eng zu offen) beim Spielverhalten fliessend, oder gibt es einen sichtbaren Bruch (Zögern, Pause, sichtbare Neuorientierung)? \hfill (Spielfluss)
\end{enumerate}

\subsubsection{Semistrukturiertes Interview}

Das Interview findet im Anschluss an das Post-Questionnaire statt und nutzt die Videoaufnahme des Playtests, um gezielt auf die während der Beobachtung notierten Auffälligkeiten einzugehen. Die folgenden Fragen dienen als Leitfaden und werden je nach den tatsächlich beobachteten Auffälligkeiten der jeweiligen Testperson angepasst und ergänzt.

\begin{enumerate}[label=\textbf{I\arabic*}, leftmargin=1.5cm, itemsep=0.5em]
\item Du bist manchmal einfach stehen geblieben und nix gemacht nach einem Raum: Hab die map angesehen, den nächsten raum, die musik und nach bugs gesucht
\item Du hast hier [Reaktion/Aussage] gezeigt bzw. gesagt. Kannst du beschreiben, warum?
\item Woran haben deine tode meistens gelegen du bist ja 3 mal gestorben: Ich weiss nicht passier halt, ich hab manchmal overcommited und habe hits kassiert und habe dann nicht gepunished.
\item Wie fühlen sich kleine räume anders and als grosse: kleine sind übersichtlicher, es gibt aber die gefahr in eine ecke gedrengt zu werden. Die grossen haben halt mehr gegner was sie für mich schwerer macht.
\item Bogenschützen zu erst und danach den rest. und hab gleichzeitig versucht zu groupen
\item tode waren sehr punishing weil du alles neu machen musstest aber ich empfand sie nicht als unfair, weil ich sie kommen gesehen habe. manchmal konnte ich aber nicht immer ganz einschäten wie viel damage ich kassiere wenn ich jetzt eine gruppe angreife
\item Hat hauptsächlich damit zu tun gehabt wie viel leben ich habe. wenig leben = stress.
\item Das design finde ich geil und den artstyle habe ich gefühlt.ich hätte gerne noch ein extra feature gehabt, dass sich ein balken aufläd für eine special attack. Fliegende gegner wären cool. Von der luft abgeschossen zu werden war sehr punishing. Ein checkpoint wäre auch noch geil gewesen. Combos waren nicht so intuitiv, weil ich mehr mixes gewohnt nicht rythmen. Das beste am Spiel war wirklich der Combat und das Movement.
\end{enumerate}

\subsection{Abbildungsverzeichnis}

\listoffigures

\subsection{KI-Deklaration}

Im Rahmen dieser Arbeit wurden KI-gestützte Werkzeuge unterstützend eingesetzt. Ihr Einsatz beschränkte sich auf die Fehlersuche im Quellcode (Debugging), die Identifikation möglicher Probleme in frühen Entwicklungsphasen sowie die sprachliche Überprüfung des Begleitberichts.

Wurde bei der Entwicklung einer Quelldatei KI oder ein anderes externes Hilfsmittel verwendet, so ist dies im Header der entsprechenden Datei dokumentiert. Sämtliche zur Bewertung eingereichten Codeabschnitte wurden von den Autoren selbst entwickelt. Auch alle Abbildungen und Grafiken wurden ohne den Einsatz generativer KI erstellt.

KI-Systeme wurden nicht zur Erstellung inhaltlicher Bestandteile dieser Arbeit verwendet. Insbesondere stammen die Forschungsfrage, die Analysen, die technischen Konzepte, die Schlussfolgerungen sowie die schriftlichen Inhalte des Begleitberichts von den Autoren. Der Einsatz von KI beschränkte sich auf sprachliche Korrekturen, formale Qualitätskontrollen sowie die Überprüfung, ob die im Projektvertrag definierten Anforderungen eingehalten wurden.

\end{document}
